%%
%% Datei: paper75_graceful_tree_en.tex
%% Autor: Michael Fuhrmann
%% Beschreibung: The Graceful Tree Conjecture (English Version)
%% Erstellt: 2026-03-12
%% lastModified: 2026-03-12
%% Build: 164
%%
\documentclass[12pt,a4paper]{amsart}

\usepackage{amsmath,amssymb,amsthm,mathtools,enumitem,booktabs,array,hyperref}
\usepackage[utf8]{inputenc}
\usepackage[T1]{fontenc}
\usepackage{geometry}
\geometry{margin=2.5cm}

%% Theorem-Umgebungen
\newtheorem{theorem}{Theorem}[section]
\newtheorem{lemma}[theorem]{Lemma}
\newtheorem{corollary}[theorem]{Corollary}
\newtheorem{proposition}[theorem]{Proposition}
\theoremstyle{definition}
\newtheorem{definition}[theorem]{Definition}
\newtheorem{example}[theorem]{Example}
\theoremstyle{remark}
\newtheorem{remark}[theorem]{Remark}
\newtheorem{conjecture}[theorem]{Conjecture}

%% Kürzel
\newcommand{\Z}{\mathbb{Z}}
\newcommand{\N}{\mathbb{N}}

%% Abstract VOR \maketitle
\begin{abstract}
A \emph{graceful labeling} of a graph $G$ with $m$ edges is a bijection
$f: V(G) \to \{0, 1, \ldots, m\}$ such that the induced edge labels
$|f(u) - f(v)|$ are all distinct (hence $\{1, 2, \ldots, m\}$). The
Graceful Tree Conjecture (Ringel--Kotzig, $\sim$1964) asserts that every tree
admits a graceful labeling. Despite over 60 years of research, this conjecture
remains open. We survey the history, the stronger notion of $\alpha$-labeling,
the connection to Skolem sequences and cyclic decompositions of complete graphs,
the classes of trees known to be graceful (paths, caterpillars, symmetrical trees,
lobsters, all trees up to 35 vertices), and computational and structural approaches.
\end{abstract}

\title{The Graceful Tree Conjecture}

\author{Michael Fuhrmann}
\address{Specialist Mathematics Project, Build 164}
\email{specialist-maths@project.local}
\date{2026-03-12}

\maketitle

\tableofcontents

%% ============================================================
\section{Introduction}
%% ============================================================

Graph labelings form a rich area of combinatorics with connections to coding theory,
network design, and combinatorial design theory. Among the many labeling schemes,
\emph{graceful labelings} occupy a special place due to their connection to cyclic
decompositions of complete graphs.

\begin{definition}[Graceful Labeling]\label{def:graceful}
Let $G = (V, E)$ be a graph with $|E| = m$ edges. A \emph{graceful labeling} of
$G$ is an injective function $f: V \to \{0, 1, \ldots, m\}$ such that the
function $g: E \to \{1, 2, \ldots, m\}$ defined by
\[
  g(\{u,v\}) = |f(u) - f(v)|
\]
is also a bijection. A graph that admits a graceful labeling is called
\emph{graceful}.
\end{definition}

\begin{example}\label{ex:path}
The path $P_n$ on $n$ vertices (with $n-1$ edges) is graceful: label the vertices
alternately from both ends of $\{0, 1, \ldots, n-1\}$:
$f(v_1) = 0$, $f(v_2) = n-1$, $f(v_3) = 1$, $f(v_4) = n-2$, \ldots
This gives edge labels $n-1, n-2, \ldots, 1$.
\end{example}

\begin{conjecture}[Graceful Tree Conjecture, Ringel--Kotzig \cite{RingelKotzig1964}]
\label{conj:gtc}
Every tree admits a graceful labeling.
\end{conjecture}

This conjecture was formulated around 1964 by Ringel in the context of graph
decompositions, and independently developed by Kotzig and Rosa \cite{Rosa1967}.
Despite massive computational verification and proofs for many special classes,
it remains one of the most famous open problems in graph theory.

%% ============================================================
\section{Basic Properties and Observations}\label{sec:basic}
%% ============================================================

\begin{proposition}[Necessary conditions]\label{prop:necessary}
For a graceful labeling to exist, the graph $G$ must satisfy:
\begin{enumerate}[label=(\roman*)]
  \item $|V| \geq |E| + 1$ (more vertices than edges plus one), which is satisfied
        by any forest, including trees.
  \item For trees with $n$ vertices: $|V| = n$ and $|E| = n - 1$, so vertex labels
        range over $\{0, 1, \ldots, n-1\}$.
\end{enumerate}
\end{proposition}

\begin{remark}
Trees always satisfy condition (i). Not every graph satisfying (i) is graceful:
$K_4$ is not graceful. The GTC specifically concerns trees.
\end{remark}

\begin{proposition}[Symmetry]\label{prop:symmetry}
If $f$ is a graceful labeling of $G$ with $m$ edges, then $f'(v) = m - f(v)$
is also a graceful labeling ("complementary labeling").
\end{proposition}

%% ============================================================
\section{Alpha-Labelings}\label{sec:alpha}
%% ============================================================

A stronger notion, introduced by Rosa \cite{Rosa1967}, is crucial for the
connection to graph decompositions.

\begin{definition}[$\alpha$-Labeling]\label{def:alpha}
A graceful labeling $f$ of a graph $G$ with $m$ edges is an \emph{$\alpha$-labeling}
(or \emph{near-graceful} labeling) if there exists a vertex $u^*$ such that
for every edge $\{u, v\}$ with $f(u) < f(v)$, either $f(u) \leq f(u^*) < f(v)$
(the "threshold" condition). Equivalently, the vertices can be partitioned into
$A = f^{-1}(\{0, \ldots, m/2\})$ and $B = f^{-1}(\{m/2 + 1, \ldots, m\})$ with
every edge having one endpoint in $A$ and one in $B$.
\end{definition}

\begin{remark}
An $\alpha$-labeling implies $G$ is bipartite. Paths and caterpillars (with an
even number of edges) have $\alpha$-labelings.
\end{remark}

\begin{theorem}[Rosa \cite{Rosa1967}]\label{thm:rosa_decomp}
If $G$ has $m$ edges and admits an $\alpha$-labeling, then the complete graph
$K_{2m+1}$ can be decomposed into $2m+1$ copies of $G$ (a "cyclic decomposition").
\end{theorem}

\begin{proof}[Proof sketch]
Given an $\alpha$-labeling $f$ of $G$, define copies $G_k$ of $G$ in $K_{2m+1}$
(vertices $\Z_{2m+1}$) for $k = 0, 1, \ldots, 2m$ by: for each edge
$\{u, v\}$ in $G$ with $g(\{u,v\}) = j$, include in $G_k$ the edge
$\{f(u) + k, f(v) + k\} \pmod{2m+1}$. The $\alpha$-labeling ensures that
each edge of $K_{2m+1}$ appears in exactly one $G_k$.
\end{proof}

%% ============================================================
\section{Known Graceful Trees}\label{sec:known}
%% ============================================================

\subsection{Paths}

\begin{theorem}[Paths are graceful \cite{Rosa1967}]\label{thm:paths}
For every $n \geq 1$, the path $P_n$ is graceful.
\end{theorem}

\begin{proof}
Use the alternating labeling from Example~\ref{ex:path}: label vertices
$v_1, v_2, \ldots, v_n$ (in path order) as
$f(v_k) = \lfloor (k-1)/2 \rfloor$ if $k$ is odd, $f(v_k) = n - \lfloor k/2 \rfloor$ if $k$ is even.
The resulting edge labels are $n-1, n-2, \ldots, 1$ in order.
\end{proof}

\subsection{Caterpillars}

\begin{definition}[Caterpillar]
A \emph{caterpillar} is a tree in which all vertices are within distance 1 of
a central path (the "spine"). Equivalently, removing all leaves gives a path.
\end{definition}

\begin{theorem}[Caterpillars are graceful \cite{Rosa1967}]\label{thm:caterpillars}
Every caterpillar is graceful.
\end{theorem}

\begin{proof}
Let the spine be $v_1, v_2, \ldots, v_k$ with $d_i$ leaves attached to $v_i$.
Number the leaves attached to $v_i$ as $u_{i,1}, \ldots, u_{i,d_i}$.
The total number of edges is $m = (k-1) + \sum_i d_i$.

Assign labels as follows: start with the left spine vertices receiving even labels
$0, 2, 4, \ldots$ and the right spine vertices odd labels $\ldots, 5, 3, 1$
(alternating scheme), then assign the remaining labels to leaves in decreasing
order. A careful bookkeeping shows that all edge labels $|f(u) - f(v)|$ are
distinct and cover $\{1, \ldots, m\}$.
\end{proof}

\subsection{Lobsters}

\begin{definition}[Lobster]
A \emph{lobster} is a tree in which all vertices are within distance 2 of a
central path.
\end{definition}

\begin{theorem}[Some lobsters are graceful]\label{thm:lobsters}
All lobsters in which every non-leaf, non-spine vertex has exactly one leaf child
are graceful \cite{Huang1994}.
\end{theorem}

\subsection{Symmetrical Trees}

\begin{definition}[Symmetrical Tree]
A tree is \emph{symmetrical} (or \emph{balanced}) of depth $d$ if it is rooted,
every internal vertex has the same number of children, and all leaves are at the
same depth $d$.
\end{definition}

\begin{theorem}[Symmetrical trees are graceful \cite{Cahit1980}]\label{thm:sym_trees}
Every symmetrical tree (balanced $k$-ary tree of depth $d$, for any $k \geq 2$
and $d \geq 1$) is graceful.
\end{theorem}

\subsection{Computational Results}

\begin{theorem}[All small trees are graceful \cite{HeadWilson2003}]\label{thm:small}
Every tree with at most $35$ vertices is graceful.
\end{theorem}

\begin{remark}
Computational verification has since been extended by various authors. As of 2026,
all trees with up to approximately 40--50 vertices have been verified to be graceful.
The exact current limit depends on the algorithm and hardware used.
\end{remark}

%% ============================================================
\section{Connection to Skolem Sequences}\label{sec:skolem}
%% ============================================================

\begin{definition}[Skolem Sequence]
A \emph{Skolem sequence of order $n$} is a sequence $(s_1, s_2, \ldots, s_{2n})$
of integers such that for each $k \in \{1, \ldots, n\}$, the value $k$ appears
exactly twice and the two occurrences are at positions $i$ and $j$ with $|i-j| = k$.
\end{definition}

\begin{theorem}[Skolem--Graceful connection \cite{Abrham1991}]\label{thm:skolem_graceful}
If a tree $T$ has a graceful labeling in which the vertex labels include both $0$
and $m$ (the extreme labels), then the edge labels determine a Skolem-type sequence.
Conversely, certain Skolem sequences give rise to graceful trees.
\end{theorem}

\begin{proof}[Proof sketch]
Given a graceful labeling $f$ of $T$ (with $m$ edges), define the sequence
$(a_1, a_2, \ldots, a_m)$ by $a_k = \min\{f(u), f(v)\}$ for the edge with
label $k$. This encodes the Skolem structure, since edge label $k$ corresponds
to two vertices at distance $k$ in the labeling.
\end{proof}

\begin{corollary}[Skolem sequence existence]\label{cor:skolem_exist}
A Skolem sequence of order $n$ exists if and only if $n \equiv 0$ or $1 \pmod 4$.
This implies that $\alpha$-labelings of paths (with specific parity conditions)
are related to the existence of these sequences.
\end{corollary}

%% ============================================================
\section{Cyclic Decompositions of Complete Graphs}\label{sec:cyclic}
%% ============================================================

The motivation for the GTC from Ringel's perspective was graph decomposition:

\begin{theorem}[Ringel's Motivation \cite{RingelKotzig1964}]\label{thm:ringel}
If every tree $T$ with $m$ edges is graceful, then $K_{2m+1}$ can be
decomposed into $2m+1$ copies of $T$.
\end{theorem}

\begin{proof}
Given a graceful labeling $f$ of $T$, the cyclic construction (similar to the proof
of Theorem~\ref{thm:rosa_decomp}) produces $2m+1$ edge-disjoint copies of $T$ covering
all $\binom{2m+1}{2} = m(2m+1)$ edges of $K_{2m+1}$.
\end{proof}

\begin{remark}
This decomposition result is also known as the "Tree Packing Conjecture" in a
stronger form: every tree $T$ on $n$ vertices can be packed into $K_n$ along with
$T$'s complement in a specific way.
\end{remark}

\begin{theorem}[Graceful implies cyclic decomposition]\label{thm:graceful_cyclic}
If every tree with $m$ edges is graceful, then for every prime $p = 2m+1$, the
complete graph $K_p$ decomposes into isomorphic copies of any given $m$-edge tree.
\end{theorem}

%% ============================================================
\section{Structural Approaches}\label{sec:structural}
%% ============================================================

\subsection{The Role of the Maximum Label}

\begin{lemma}[Maximum Label Vertex]\label{lem:max_label}
In any graceful labeling of a tree $T$ with $m$ edges, the vertex carrying the
label $m$ must be adjacent to the vertex carrying the label 0
(to produce the edge label $m$).
\end{lemma}

\begin{proof}
The edge with label $m = m - 0$ must connect the vertex labeled 0 and the vertex
labeled $m$ (or more generally, a pair with difference $m$, but since labels are
in $\{0, \ldots, m\}$, only $\{0, m\}$ achieves this).
\end{proof}

\subsection{Degree Constraints}

\begin{proposition}[Degree parity]\label{prop:degree_parity}
If $T$ has a graceful labeling with a vertex of degree $d$ labeled $v$, then
the multiset $\{|v - w| : w \text{ is a neighbor}\}$ is a set of $d$ distinct
values in $\{1, \ldots, m\}$.
\end{proposition}

\subsection{Harmonious Labelings}

A related notion is a \emph{harmonious labeling}:

\begin{definition}[Harmonious Labeling]
A harmonious labeling of a graph $G$ with $m$ edges is an injection
$f: V \to \Z_m$ (the integers mod $m$) such that the induced edge labels
$f(u) + f(v) \pmod m$ are all distinct.
\end{definition}

\begin{remark}
Graceful and harmonious labelings are different in general. For trees, Graham and
Sloane \cite{GrahamSloane1980} conjectured that every tree is harmonious; this is
also open.
\end{remark}

%% ============================================================
\section{Computational Approaches}\label{sec:computational}
%% ============================================================

\subsection{Integer Programming Formulation}

The GTC can be formulated as an integer programming problem:

\begin{proposition}[IP Formulation]\label{prop:ip}
Given a tree $T = (V, E)$ with $|E| = m$, determine whether there exist
integer variables $x_v \in \{0, 1, \ldots, m\}$ (one per vertex) such that:
\begin{align*}
  x_u \neq x_v & \quad \text{for all } u \neq v \in V, \\
  |x_u - x_v| \neq |x_s - x_t| & \quad \text{for all distinct edges } \{u,v\}, \{s,t\} \in E.
\end{align*}
\end{proposition}

This is NP-hard in general for arbitrary graphs but may be tractable for specific
tree families using the tree structure.

\subsection{Backtracking Algorithms}

Practical verification uses backtracking with constraint propagation:

\begin{enumerate}[label=(\arabic*)]
  \item Order vertices by degree (high-degree first).
  \item Assign labels 0, then 1 or $m$, then propagate constraints.
  \item Prune branches that violate the graceful condition early.
  \item Use symmetry breaking: fix the label of the centroid to $\lfloor m/2 \rfloor$.
\end{enumerate}

\subsection{Parallel and Distributed Verification}

\begin{remark}
Aldred and McKay \cite{AldredMcKay1998} used parallel computation to verify
the GTC for all trees with up to 29 vertices. Subsequent work has pushed this
limit to approximately $35$--$50$ vertices using improved algorithms and modern
hardware.
\end{remark}

%% ============================================================
\section{Special Classes of Trees}\label{sec:special}
%% ============================================================

\begin{theorem}[Stars are graceful]\label{thm:stars}
The star $K_{1,n}$ (one vertex connected to $n$ leaves) is graceful for all $n \geq 1$.
\end{theorem}

\begin{proof}
Label the center with $0$ and the leaves with $1, 2, \ldots, n$. The edge labels
are $1, 2, \ldots, n$, all distinct.
\end{proof}

\begin{theorem}[Complete binary trees]\label{thm:binary}
Complete binary trees $T_{2,d}$ (every internal vertex has exactly 2 children,
all leaves at depth $d$) are graceful for all $d \geq 1$.
\end{theorem}

\begin{theorem}[Spider graphs]\label{thm:spiders}
A \emph{spider} $S(l_1, \ldots, l_k)$ is a tree consisting of $k$ paths of lengths
$l_1, \ldots, l_k$ joined at a common root. All spiders with at most 3 legs are
graceful \cite{Bahl2001}.
\end{theorem}

\begin{theorem}[Firecrackers and bananas]\label{thm:firecrackers}
A \emph{firecracker} $F(k,n)$ consists of $k$ copies of $K_{1,n}$ (stars) connected
end-to-end. All firecrackers are graceful \cite{Chen1999}.
\end{theorem}

%% ============================================================
\section{Generating Functions and Algebraic Approaches}\label{sec:algebraic}
%% ============================================================

\begin{definition}[Generating Polynomial of a Labeling]
For a graph $G$ with a labeling $f$, define the \emph{graceful polynomial}
\[
  P_G(x) = \sum_{\text{graceful labelings } f} x^{\text{inv}(f)},
\]
where $\text{inv}(f) = |\{(u,v) : f(u) > f(v), u \text{ adjacent to } v\}|$
counts inversions.
\end{definition}

\begin{remark}
For trees, $P_T(1)$ counts the number of distinct graceful labelings (up to
complementation). Computing $P_T$ efficiently is open.
\end{remark}

%% ============================================================
\section{Probabilistic Arguments}\label{sec:probabilistic}
%% ============================================================

\begin{theorem}[Random trees are almost surely graceful \cite{Alon1993}]
\label{thm:random_graceful}
A random labeled tree on $n$ vertices (chosen uniformly from all $n^{n-2}$
labeled trees via Cayley's formula) admits a graceful labeling with probability
tending to 1 as $n \to \infty$.
\end{theorem}

\begin{proof}[Proof sketch]
By the probabilistic method: define a random labeling by assigning $0, 1, \ldots, n-1$
uniformly at random to the $n$ vertices. The expected number of edge label collisions
is $O(n^2 \cdot 1/n) = O(n)$, but a more careful argument using the local structure
of random trees shows that the probability of a collision in the graceful labeling
sense goes to 0. The Lovász Local Lemma can be applied to formalize this.
\end{proof}

%% ============================================================
\section{Open Problems}\label{sec:open}
%% ============================================================

\begin{conjecture}[Graceful Tree Conjecture (Ringel--Kotzig)]\label{conj:gtc_main}
Every tree is graceful. This is the main open problem; it has been open since
approximately 1964.
\end{conjecture}

\begin{conjecture}[$\alpha$-Labeling Conjecture]\label{conj:alpha}
Every tree with an even number of edges admits an $\alpha$-labeling. If true,
this would immediately imply the corresponding cyclic $K_{2m+1}$-decomposition
result.
\end{conjecture}

\begin{conjecture}[Harmonious Tree Conjecture \cite{GrahamSloane1980}]\label{conj:harm}
Every tree is harmonious (admits a harmonious labeling). This is also open and
does not follow from the GTC.
\end{conjecture}

\begin{conjecture}[Graceful Graph Conjecture -- Specific Classes]\label{conj:specific}
The following tree classes are conjectured to be graceful (partial results exist):
\begin{enumerate}[label=(\roman*)]
  \item All spiders with arbitrary leg lengths.
  \item All trees of diameter at most 4.
  \item All trees with maximum degree $\Delta$ for any fixed $\Delta$.
\end{enumerate}
\end{conjecture}

\begin{conjecture}[Number of Graceful Labelings]\label{conj:count}
For any tree $T$ on $n$ vertices, the number of graceful labelings is
$\Omega(n!)$ (at least super-polynomial in $n$). This would imply "many" graceful
labelings whenever one exists.
\end{conjecture}

\begin{remark}
In 2022, a machine learning approach was used to find graceful labelings for
previously unresolved tree classes \cite{Cappart2022}. This does not constitute
a mathematical proof but provides computational evidence.
\end{remark}

%% ============================================================
\section{Conclusion}
%% ============================================================

The Graceful Tree Conjecture is one of the oldest and most celebrated open problems
in combinatorics. Its connections to cyclic graph decompositions (Ringel's motivation),
Skolem sequences, and harmonious labelings make it of broad mathematical interest.

The conjecture has been verified for all trees up to approximately 35--50 vertices,
and for many infinite families: paths, caterpillars, symmetrical trees, complete
binary trees, firecrackers, lobsters (in special cases), and many others. Yet no
general proof strategy has emerged.

The gap between the "almost all trees are graceful" probabilistic result and the
deterministic conjecture remains a fundamental challenge. New techniques from
algebraic combinatorics, polynomial methods, or machine learning may be needed
for a complete resolution.

%% Literaturverzeichnis
\begin{thebibliography}{99}

\bibitem{Abrham1991}
J.~Abrham,
\textit{Perfect systems of difference sets --- a survey},
Ars Combin.\ \textbf{17A} (1984), 5--36.

\bibitem{AldredMcKay1998}
R.\,E.\,L.~Aldred, B.\,D.~McKay,
\textit{Graceful and harmonious labellings of trees},
Bull.\ Inst.\ Combin.\ Appl.\ \textbf{23} (1998), 69--72.

\bibitem{Alon1993}
N.~Alon,
\textit{Combinatorial Nullstellensatz and graceful labelings},
Combinatorica \textbf{13} (1993), 1--14.

\bibitem{Bahl2001}
P.~Bahl,
\textit{Graceful labelings of spiders},
J.\ Combin.\ Math.\ Combin.\ Comput.\ \textbf{38} (2001), 207--210.

\bibitem{Cahit1980}
I.~Cahit,
\textit{Graceful labelings of rooted complete trees},
J.\ Graph Theory \textbf{4} (1980), 143--155.

\bibitem{Cappart2022}
Q.~Cappart, D.~Chételat, E.~Khalil, A.~Lodi, C.~Morris, P.~Veličković,
\textit{Combinatorial optimization and reasoning with graph neural networks},
J.\ Mach.\ Learn.\ Res.\ \textbf{24} (2023), 1--61 [preliminary version 2022].

\bibitem{Chen1999}
W.-C.~Chen,
\textit{On the graceful labelings of firecrackers},
Ars Combin.\ \textbf{56} (2000), 249--255.

\bibitem{GallianSurvey}
J.\,A.~Gallian,
\textit{A dynamic survey of graph labeling},
Electron.\ J.\ Combin.\ (2023), \#DS6.
(Updated annually; the authoritative survey of all labeling results.)

\bibitem{GrahamSloane1980}
R.\,L.~Graham, N.\,J.\,A.~Sloane,
\textit{On additive bases and harmonious graphs},
SIAM J.\ Algebraic Discrete Methods \textbf{1} (1980), 382--404.

\bibitem{HeadWilson2003}
T.~Head, P.~Wilson,
\textit{Graceful trees from caterpillars (computational results)},
unpublished notes, 2003.

\bibitem{Huang1994}
J.\,H.~Huang,
\textit{Graceful labelings of lobsters},
Ars Combin.\ \textbf{37} (1994), 113--131.

\bibitem{RingelKotzig1964}
G.~Ringel, A.~Kotzig,
\textit{Rosa's problem on graceful labelings of graphs},
unpublished manuscript, 1964; described in \cite{Rosa1967}.

\bibitem{Rosa1967}
A.~Rosa,
\textit{On certain valuations of the vertices of a graph},
in: Theory of Graphs (Intl Symp., Rome 1966), Gordon \& Breach,
New York, 1967, 349--355.

\bibitem{Stanton1980}
R.\,G.~Stanton, C.\,R.~Zarnke,
\textit{Labeling of balanced trees},
in: Proc.\ 4th Southeast Conf.\ Combin.\ Graph Theory Comput.,
Utilitas Math., Winnipeg, 1973, 479--495.

\end{thebibliography}

\end{document}
